% ============================================================
\section{Methodology}
\label{sec:methodology}
% ============================================================

\subsection{Overview}
\label{sec:method-overview}
\lipsum[3]

% Example figure (put files in assets/figures/ and reference by name only):
% \begin{figure}[t]
%   \centering
%   \includegraphics[width=0.7\linewidth]{example-figure}
%   \caption{Overview of the proposed approach.}
%   \label{fig:overview}
% \end{figure}

\subsection{Problem Formulation}
\label{sec:formulation}
Let $x \in \R^d$ denote an input. We seek
\begin{equation}
  \theta^\star = \argmin_{\theta} \; \E_{(x,y)}\big[\ell(f_\theta(x), y)\big].
  \label{eq:objective}
\end{equation}
\Cref{eq:objective} defines the objective.

\subsection{Data}
\label{sec:data}
\todo{Describe data sources and preprocessing.}

% Example table:
\begin{table}[t]
  \centering
  \caption{Example table using \texttt{booktabs}.}
  \label{tab:example}
  \begin{tabular}{lcc}
    \toprule
    Method & Metric A & Metric B \\
    \midrule
    Baseline & 0.00 & 0.00 \\
    Ours     & 0.00 & 0.00 \\
    \bottomrule
  \end{tabular}
\end{table}

\subsection{Evaluation}
\label{sec:evaluation}
\todo{Describe metrics and experimental setup (see \cref{tab:example}).}
